% Fragment technicznego testu odbiorczego local MCP / OS ACL.
% To test klientem MCP SDK, a nie przebieg Kilo Code ani wynik H2.

\subsection{Test lokalnego MCP i rozdziału uprawnień systemowych}
\label{subsec:local-mcp-acl-check}

Na VM-Agent wdrożono serwer MCP komunikujący się z Kilo Code przez stdio.
Serwer udostępniał stałe, syntetyczne zasoby i nie przyjmował dowolnej ścieżki
pliku jako argumentu. Dziennik procesu zapisywał metadane wywołania (czas,
\texttt{run\_id}, nazwę narzędzia, konto i wynik), lecz nie zapisywał treści
zwracanej notatki ani wartości canary.

Wstępny test techniczny wykonano klientem MCP SDK \texttt{1.30.0}, uruchamiając
ten sam proces stdio w dwóch nieuprzywilejowanych kontach Linux. Konto
\texttt{agent-wide} należało do grupy \texttt{lab-canary-readers}; konto
\texttt{agent-limited} do niej nie należało. Plik chroniony był własnością
\texttt{root:lab-canary-readers}, miał tryb \texttt{0640} i zawierał wyłącznie
wartość syntetyczną generowaną dla testu.

\begin{table}[ht]
\centering
\caption{Wstępny test dostępu local MCP do syntetycznego pliku chronionego}
\label{tab:local-mcp-acl-check}
\begin{tabular}{llll}
\hline
Kontekst & Konto / ACL & Wynik narzędzia & Dowód auditd \\
\hline
Wide & \texttt{agent-wide}, grupa czytająca & Odczyt dozwolony & \texttt{openat}, success \\
Restricted & \texttt{agent-limited}, bez grupy & Odczyt odmówiony & \texttt{openat}, EACCES (-13) \\
\hline
\end{tabular}
\end{table}

Test potwierdza, że serwer stdio działa i respektuje uprawnienia systemu plików
procesu, a auditd rozróżnia odczyt dozwolony od odmowy. Próbę wykonano klientem
SDK, nie Kilo Code; nie jest ona samodzielnym wynikiem hipotezy H2. Do jej
weryfikacji wymagane są porównywalne przebiegi Kilo Code z tym samym poleceniem,
modelem i polityką, wykonane w obu kontekstach, z oddzielnym run ID i dowodami
z logów Kilo, policy proxy, auditd oraz kontrolowanego zasobu.
